\section{Cơ sở lý thuyết của DFS}

DFS là thuật toán duyệt đồ thị theo nguyên tắc đi sâu theo từng nhánh trước, rồi mới quay lui khi không thể đi tiếp.

Trong bài toán flood fill, DFS được mô tả qua các yếu tố sau:

\textbf{Cấu trúc dữ liệu:}  
Thuật toán sử dụng một stack. Đỉnh được đưa vào sau sẽ được xử lý trước.

\textbf{Luật chuyển trạng thái:}  
Một ô chỉ được đưa vào stack nếu nó hợp lệ: chưa mở, không bị đánh dấu, và không phải là mìn. Điều này giúp tránh lặp lại và đảm bảo thuật toán không mở sai ô.

\textbf{Điều kiện dừng:}  
Khi DFS gặp một ô có số lớn hơn 0 (tức là ô nằm ở viền vùng trống), thuật toán vẫn mở ô đó nhưng không lan tiếp sang các ô lân cận. Đây chính là cơ chế “cắt nhánh” giúp flood fill dừng đúng chỗ.
